\documentclass{beamer}
\usepackage[utf8]{inputenc}

\usetheme{Madrid}
\usecolortheme{default}
\usepackage{amsmath,amssymb,amsfonts,amsthm}
\usepackage{mathtools}
\usepackage{txfonts}
\usepackage{tkz-euclide}
\usepackage{listings}
\usepackage{adjustbox}
\usepackage{array}
\usepackage{gensymb}
\usepackage{tabularx}
\usepackage{gvv}
\usepackage{lmodern}
\usepackage{circuitikz}
\usepackage{tikz}
\lstset{literate={·}{{$\cdot$}}1 {λ}{{$\lambda$}}1 {→}{{$\to$}}1}
\usepackage{graphicx}

\setbeamertemplate{page number in head/foot}[totalframenumber]

\usepackage{tcolorbox}
\tcbuselibrary{minted,breakable,xparse,skins}



\definecolor{bg}{gray}{0.95}
\DeclareTCBListing{mintedbox}{O{}m!O{}}{%
  breakable=true,
  listing engine=minted,
  listing only,
  minted language=#2,
  minted style=default,
  minted options={%
    linenos,
    gobble=0,
    breaklines=true,
    breakafter=,,
    fontsize=\small,
    numbersep=8pt,
    #1},
  boxsep=0pt,
  left skip=0pt,
  right skip=0pt,
  left=25pt,
  right=0pt,
  top=3pt,
  bottom=3pt,
  arc=5pt,
  leftrule=0pt,
  rightrule=0pt,
  bottomrule=2pt,
  toprule=2pt,
  colback=bg,
  colframe=orange!70,
  enhanced,
  overlay={%
    \begin{tcbclipinterior}
    \fill[orange!20!white] (frame.south west) rectangle ([xshift=20pt]frame.north west);
    \end{tcbclipinterior}},
  #3,
}
\lstset{
    language=C,
    basicstyle=\ttfamily\small,
    keywordstyle=\color{blue},
    stringstyle=\color{orange},
    commentstyle=\color{green!60!black},
    numbers=left,
    numberstyle=\tiny\color{gray},
    breaklines=true,
    showstringspaces=false,
}

\numberwithin{equation}{section}
\lstset{
  language=Python,
  basicstyle=\ttfamily\small,
  keywordstyle=\color{blue},
  stringstyle=\color{orange},
  numbers=left,
  numberstyle=\tiny\color{gray},
  breaklines=true,
  showstringspaces=false
}


\title{Problem 5.7.3}
\author{Sarvesh Tamgade}

\date{\today} 
\begin{document}

\begin{frame}
\titlepage
\end{frame}

\section{Question}
\begin{frame}{Question}
\textbf{Question}:
For the matrix
\[
\myvec{
1 & 1 & 1 \\
1 & 2 & -2 \\
2 & -1 & 3
},
\]
show that
\begin{align}
A^3 - 6A^2 + 5A + 11I &= 0.
\end{align}
Find \(A^{-1}\). 
\end{frame}

\section{Solution}
\begin{frame}[fragile]
    \frametitle{Solution}
\textbf{Solution:} \
The characteristic equation is
\begin{align}
\left|\myvec{
1-\lambda & 1 & 1 \\
1 & 2-\lambda & -2 \\
2 & -1 & 3-\lambda
} \right| &= 0,
\end{align}
which expands to
\begin{align}
\lambda^3 - 6 \lambda^2 + 5 \lambda + 11 &= 0.
\end{align}
By Cayley-Hamilton theorem,
\begin{align}
A^3 - 6A^2 + 5A + 11I &= 0.
\end{align}
Rearranging,
\begin{align}
11I &= -A^3 + 6A^2 - 5A,
\end{align}

\end{frame}



\begin{frame}[fragile]
    \frametitle{Solution}

and multiplying both sides by \(A^{-1}\),
\begin{align}
A^{-1} &= \frac{1}{11}(-A^2 + 6A - 5I) \\
&= \frac{1}{11} \myvec{
-3 & 4 & 4 \\
7 & 0 & -3 \\
5 & -3 & 0
}.
\end{align}

\textbf{Final Answer:}
\begin{align}
A^{-1} &= \myvec{
-\frac{3}{11} & \frac{4}{11} & \frac{4}{11} \\
\frac{7}{11} & 0 & -\frac{3}{11} \\
\frac{5}{11} & -\frac{3}{11} & 0
}.
\end{align}
\end{frame}
\section{ C Code}
\begin{frame}[fragile]
\frametitle{C Code }
\begin{lstlisting}[language=C]
#include <stdio.h>
#include "matfun.h"

void matmult(double A[SIZE][SIZE], double B[SIZE][SIZE], double res[SIZE][SIZE]) {
    for(int i=0; i<SIZE; i++) {
        for(int j=0; j<SIZE; j++) {
            res[i][j] = 0.0;
            for(int k=0; k<SIZE; k++) {
                res[i][j] += A[i][k] * B[k][j];
            }
        }
    }
}

\end{lstlisting}
\end{frame}
\begin{frame}[fragile]
\frametitle{C Code }
\begin{lstlisting}[language=C]
void matscalarmult(double scalar, double A[SIZE][SIZE], double res[SIZE][SIZE]) {
    for(int i=0; i<SIZE; i++) {
        for(int j=0; j<SIZE; j++) {
            res[i][j] = scalar * A[i][j];
        }
    }
}

void matadd(double A[SIZE][SIZE], double B[SIZE][SIZE], double res[SIZE][SIZE]) {
    for(int i=0; i<SIZE; i++) {
        for(int j=0; j<SIZE; j++) {
            res[i][j] = A[i][j] + B[i][j];
        }
    }
}
\end{lstlisting}
\end{frame}
\begin{frame}[fragile]
\frametitle{C Code }
\begin{lstlisting}[language=C]
void matsub(double A[SIZE][SIZE], double B[SIZE][SIZE], double res[SIZE][SIZE]) {
    for(int i=0; i<SIZE; i++) {
        for(int j=0; j<SIZE; j++) {
            res[i][j] = A[i][j] - B[i][j];
        }
    }
}

void matidentity(double I[SIZE][SIZE]) {
    for(int i=0; i<SIZE; i++) {
        for(int j=0; j<SIZE; j++) {
            I[i][j] = (i == j) ? 1.0 : 0.0;
        }
    }
}
\end{lstlisting}
\end{frame}
\begin{frame}[fragile]
\frametitle{C Code }
\begin{lstlisting}[language=C]
void compute_inverse(double A[SIZE][SIZE], double A_inv[SIZE][SIZE]) {
    double A2[SIZE][SIZE], negA2[SIZE][SIZE];
    double sixA[SIZE][SIZE], neg5I[SIZE][SIZE];
    double temp[SIZE][SIZE], I[SIZE][SIZE];

    matmult(A, A, A2);
    matscalarmult(-1, A2, negA2);
    matscalarmult(6, A, sixA);

    matidentity(I);
    matscalarmult(-5, I, neg5I);

    matadd(negA2, sixA, temp);
    matadd(temp, neg5I, temp);

    matscalarmult(1.0/11.0, temp, A_inv);
}
\end{lstlisting}
\end{frame}
\begin{frame}[fragile]
\frametitle{Python Code }
\begin{lstlisting}[language=Python]
iimport ctypes
import numpy as np

SIZE = 3

lib = ctypes.CDLL('./libmatfun.so')

Matrix3x3 = ctypes.c_double * (SIZE*SIZE)

def np_to_c_matrix(np_mat):
    flat = np_mat.flatten()
    return Matrix3x3(*flat)
\end{lstlisting}

\end{frame}
\begin{frame}[fragile]
    \frametitle{Python Code}
    \begin{lstlisting}
def c_matrix_to_np(c_mat):
    return np.array(list(c_mat)).reshape((SIZE, SIZE))

# Prototype for compute_inverse
lib.compute_inverse.argtypes = [ctypes.POINTER(ctypes.c_double), ctypes.POINTER(ctypes.c_double)]

# Matrix A
A_np = np.array([[1,1,1],
                [1,2,-2],
                [2,-1,3]], dtype=np.double)

A_c = np_to_c_matrix(A_np)
A_inv_c = Matrix3x3()
\end{lstlisting}
\end{frame}

\begin{frame}[fragile]
    \frametitle{Python Code }
    \begin{lstlisting}
# Call the C function
lib.compute_inverse(A_c, A_inv_c)

A_inv_np = c_matrix_to_np(A_inv_c)

print("Inverse matrix A^-1:")
for row in A_inv_np:
    print(" ".join(f"{x: .4f}" for x in row))


\end{lstlisting}
\end{frame}
\end{document}
